\documentclass[11pt]{article}
\usepackage{amsmath,amssymb,amsthm}
\usepackage{fullpage}
\usepackage{amsmath,amssymb,amsthm}
\usepackage[T1]{fontenc}
\usepackage{hyperref}
\newcommand{\M}{\mathcal M}
\newcommand{\T}{\mathcal T}
\newcommand{\cl}[1]{\overline{#1}}
\newcommand{\Mod}{\operatorname{Mod}}
\newcommand{\Aut}{\operatorname{Aut}}
\newtheorem{theorem}{Theorem}
\newtheorem{lemma}[theorem]{Lemma}
\newtheorem{proposition}[theorem]{Proposition}
\newtheorem{corollary}[theorem]{Corollary}
\theoremstyle{remark}
\newtheorem{remark}[theorem]{Remark}

\newcommand{\fptitle}{Review of Problem 9 -- ProofCouncil}
\newcommand{\reviewtag}{Reviewer 3}
\usepackage{fp-review}

\begin{document}
\maketitle

\section{Review}
The solution acknowledges a partial result, only proving cases where the boundary stratum has at least two nodes (Corollary 3) or the totally geodesic surface is nonarithmetic (Corollary 11).
The remaining case is proved only under an additional nonstandard ``lifting property'' (L) (Theorem 20).
The intended solution did not need any arithmetic input, instead relying only on the (contemporaneous) work of Arana-Herrera--Wright and Benirschke--Dozier--Rached.
(The former was cited while the latter was not.)

The proposed strategy from the introduction is unclear, and it does not represent all the directions the paper attempts.
This includes a very technical computation in Lemma 13 concerning ``sewing obstructions,'' which again seems unnecessary for the original problem.
Step 2 in Theorem 20 might be getting at the right idea for the general case, but we cannot confirm this, as it is buried among many pages of unnecessary material.

\paragraph{Correctness.}
\begin{itemize}
	\item  1: essentially flawless mathematical argument
	\item  0.5: minor errors that are easily corrected
	\item  0: blank solution, or acknowledged partial result
	\item  $-0.5$: major errors
	\item  $-1$: proposed proof makes no progress towards the result.
\end{itemize}
Note that points should be docked for wrong statements/ solutions.

\noindent Score: 0

\paragraph{Exposition.}
\begin{itemize}
	\item  1: well-written and clear solution.
	\item  0: argument that is possible to follow with serious effort
	\item  $-1$: very poor exposition
\end{itemize}

\noindent Score: $-1$

\paragraph{Attributions.}
\begin{itemize}
	\item  1: solutions contains fully appropriate attributions
		to related work and theorems being used.
	\item  0: moderate issues with missing attributions or hallucinated references
	\item  $-1$: blatant plagiarism or hallucinations
\end{itemize}

\noindent Score: 0

\newpage
\section{Submission with inline comments}


\noindent\textbf{Partial result: not a complete solution (unreviewed).}\par

\begin{center}
\textbf{Rigidity of a constant boundary factor of a deformed
totally geodesic surface: partial results}
\end{center}

\noindent\textbf{Abstract.}
Under boundary hypotheses (i)--(ii) we ask whether the constant
factor \(Y_t\) of the boundary curve of a deformed totally geodesic
surface \(V_t\subset\M_g\) is independent of \(t\). This is proved
when the boundary stratum has at least two nodes or \(V_0\) is
nonarithmetic; in the remaining one-node arithmetic case
it is proved only under an additional lifting property (L), so the
problem is \emph{not} fully resolved.

\section{Introduction}

\paragraph{The problem.}
Let \(\M_g\) be the moduli space of genus \(g\) Riemann surfaces and
\(\cl{\M}_g\) its Deligne--Mumford compactification. The boundary
\(\partial\M_g\) is partitioned into strata \(\Delta_G\) parametrizing
stable curves with dual graph \(G\); \(\Delta_{G'}\) lies
on the boundary of \(\Delta_G\) when \(G'\) is obtained from \(G\) by
pinching further curves. Throughout, closures are taken in \(\cl{\M}_g\).

Let \(V_0\subset\M_g\) be a smooth complex surface, totally geodesic
for the Teichm\"uller metric, and let \(V_t\), \(t\) in the unit disc
\(\Delta\), be a holomorphic family of subvarieties deforming it:
there is a complex submanifold \(\mathcal V\subset\M_g\times\Delta\)
with \(V_t\) its fibre over \(t\), and this family is assumed smoothly
trivial. The \(V_t\) for \(t\ne0\) are \emph{not} assumed totally
geodesic. Two hypotheses are imposed:
(i) for every boundary stratum \(\Delta_G\), the intersections
\(\cl{V_t}\cap\Delta_G\) form a smoothly trivial family;
(ii) these boundary families are compatible with \(V_0\): if
\(\Delta_{G'}\) lies on the boundary of \(\Delta_G\) and
\(U\subset\cl{\M}_g\) is open with
\(\cl{\cl{V_0}\cap\Delta_G}\cap\Delta_{G'}\cap U
   =\cl{V_0}\cap\Delta_{G'}\cap U\),
then for all \(t\),
\(\cl{\cl{V_t}\cap\Delta_G}\cap\Delta_{G'}\cap U
   =\cl{V_t}\cap\Delta_{G'}\cap U\).
Fix a dual graph \(\Gamma\) whose stratum is a product
\(\Delta_\Gamma=\M_{g_1,n_1}\times\cdots\times\M_{g_k,n_k}\), with
projections \(\pi_j\), and put \(\partial V_t=\cl{V_t}\cap\Delta_\Gamma\).
Suppose \(\dim\partial V_t=1\) and that for some \(i\),
\(\pi_i(\partial V_t)=\{Y_t\}\) is a single point for every \(t\).
\emph{The task is to prove \(Y_t=Y_0\) for all \(t\).}

For a boundary stratum \(S\) we write \(B_S(t)=\cl{V_t}\cap S\),
with its reduced structure; since \(V_t\) is algebraic,
\(\dim B_S(t)\le1\), and \(\partial V_t=B_{\Delta_\Gamma}(t)\).
We read ``subvariety'' as closed algebraic, and ``smoothly trivial
family'' in the embedded sense: the union of the fibres is a smooth
submanifold of the target times the disc, and projection to the
disc is a smoothly trivial submersion (merely diffeomorphic fibres
would be weaker).
For the boundary families in (i), holomorphicity then
\emph{follows} (Lemma~\ref{lem:regularity}). All moduli spaces are
coarse unless a lift to Teichm\"uller space is explicitly specified.

\paragraph{What is proved.}
Without further hypotheses: (1) if \(\Gamma\) has at least two
edges, \(Y_t=Y_0\) (Corollary~\ref{cor:deep}); (2) if the field of
definition \(K_N\) of the affine invariant manifold attached to a
marked lift \(N\subset\T_g\) of \(V_0\) (Section~\ref{sec:teich}) is
not \(\mathbb Q\), then \(Y_t=Y_0\)
(Corollary~\ref{cor:nonarithmetic}); (3) in the sole unresolved case,
\(\Delta_\Gamma=\M_{h,1}\times\M_{m,1}\) is a one-node separating
stratum with \(Y_t\in\M_{h,1}\), the active (moving) factor is an
arithmetic Teichm\"uller curve in \(\M_{m,1}\), \(m\ge2\) (one
generated by a square-tiled surface), whose generating quadratic
differential vanishes at the marked point
(Theorem~\ref{thm:zero}, Proposition~\ref{prop:active-field});
(4) every two-node cusp (nodal limit point) of the central (\(t=0\))
one-node boundary curve is isolated in \(\cl{V_0}\cap S\) for its
two-node stratum \(S\), and by (i)--(ii) has a unique holomorphic
continuation
\(p(t)\) (Proposition~\ref{prop:two-corner},
Corollary~\ref{cor:corner-section}).
What is \emph{not} proved is that \(p(t)\) is stationary, or that
its normalization components other than \(Y_t\) are fixed.
Theorem~\ref{thm:conditional} proves \(Y_t=Y_0\) in general under an
explicit lifting property (L) on the boundary curves
(Section~\ref{sec:lifting}); Section~\ref{sec:open} records why (L)
has not been derived from (i)--(ii).

\paragraph{Strategy.}
The mechanism is a \emph{corner argument}
(Proposition~\ref{prop:fixed-cusp}): if a stable curve \(p\) with
more nodes lies in the closure of \(\partial V_0\), is isolated in
\(\cl{V_0}\cap S\) for its own stratum \(S\), and \emph{stays} in
\(\cl{V_t}\cap S\) for small \(t\), then \(Y_t\) lies in the finite
set of smooth pointed normalization components of \(p\).
Hypothesis (ii) gives this stationarity when \(S\) has at
least three nodes (Lemma~\ref{lem:deep-points}); the hard case is a
one-node stratum, whose simplest corners are two-node curves: we
obtain their isolation (Proposition~\ref{prop:two-corner}) but not
their stationarity. Sections~\ref{sec:teich}--\ref{sec:sewing} use the
boundary theory of \cite{AHW}, cylinder deformations
\cite{WrightCylinder} and a second-order sewing expansion of period
matrices \cite{HuNorton}; Section~\ref{sec:lifting} introduces (L)
and the rigidity of Teichm\"uller curve maps \cite{McMullen}, and
Section~\ref{sec:conditional} runs the corner argument under (L).

\section{Consequences of hypotheses (i)--(ii)}\label{sec:hyp}

\begin{lemma}[Holomorphicity of the boundary families]\label{lem:regularity}
Under the above interpretation, the boundary families in (i)
are holomorphic submersions. In particular \(t\mapsto Y_t\)
is holomorphic, so local constancy at \(0\) implies constancy
on the whole disc.
\end{lemma}
\begin{proof}
\emph{Step 1: a generic projective envelope.}
Embed \(\cl{\M}_g\) in projective space and choose an affine graph
\((z,F(z,t))\) for one connected component of the interior
family \(V_t\), near any chosen parameter \(t=0\).
Every \(F_j(\,\cdot\,,t)\) is algebraic over \(\mathbb C(z)\).
For each degree \(n\), evaluate the finitely many monomials
in \((z,F_j)\) of degree at most \(n\) and take their
Taylor-coefficient matrix in \(z\).
Linear dependence of its columns is a closed condition on
\(t\), given by holomorphic minors, and these closed sets cover
the disc, so Baire's theorem and the identity theorem make all
maximal minors vanish identically for some \(n\).
Over \(K=\operatorname{Frac}\mathbb C\{t\}\), finitely many
rows span the row space; a nonzero vector in their kernel
annihilates every row and gives a polynomial relation for \(F_j\).
Thus \(K(z,F)\) is finite over \(K(z)\), and
\(K\) is algebraically closed in this field:
differentiating the minimal polynomial of any element
algebraic over \(K\) shows that all its \(z\)-derivatives
vanish, so the element is meromorphic in \(t\) alone.
The prime ideal of all relations therefore defines a
geometrically integral generic variety \cite[Tag 037P]{Stacks}
(reducedness is automatic in characteristic zero).
\review{Doing algebraic geometry at the level of the Stacks project should be completely unneccessary for this problem.}
Take its
projective closure over \(\mathbb C\{t\}\) and spread its finitely
many homogeneous equations over a finitely generated complex
subalgebra \(R\subset\mathbb C\{t\}\); the generic fibre
is geometrically integral, since this holds after extending
\(\operatorname{Frac}R\) to \(K\). Generic flatness and
generic geometric integrality
\cite[Tags 052A, 0559, 0578]{Stacks} give integral fibres
of dimension two after inverting some \(0\ne f\in R\).

The nonzero germ \(f(t)\) has no zeros on a small punctured disc,
where each envelope fibre contains the original graph and equals
the individual projective closure of that component by
irreducibility and dimension. Repeat for the finitely many continued components;
the central envelope may have extra components, none of which are
excluded or used here.

\emph{Step 2: the boundary family is complex.}
Consequently \(\bigcup_{t\ne0}B_S(t)\times\{t\}\) is analytic.
By (i) its underlying set, including \(t=0\), is an embedded
smooth family. At the regular points of the analytic part its
tangent spaces are complex; density and continuity of tangent
spaces make them complex everywhere, including at \(t=0\), and
a smooth embedded submanifold with complex tangent spaces
is a complex submanifold (its local graph maps have complex-linear
differentials, hence are holomorphic).
Its projection, already a smooth submersion, is therefore
a holomorphic submersion. This applies near any parameter.
Finally, a local holomorphic section followed by \(\pi_i\) gives
\(Y_t\), and the identity theorem proves the last assertion.
\end{proof}

The next lemmas exploit the literal quantifier over open sets in
(ii), applied to the particular open set
\(U=\cl{\M}_g\setminus\cl{B_S(0)}\), for which both sides of the
antecedent are empty.

\begin{lemma}[Fixed curves in deeper strata]\label{lem:deep}
If \(S\) has at least two nodes, every curve component of \(B_S(t)\)
is contained in \(B_S(0)\). A curve component continued from
\(B_S(0)\) is therefore fixed, as a subset, for all small \(t\).
\end{lemma}

\begin{proof}
Choose a boundary stratum \(H\) obtained by smoothing one node of
\(S\), and set \(U=\cl{\M}_g\setminus\cl{B_S(0)}\).
Both sets in the antecedent of (ii), for \(H,S,U\), are empty, since
\(\cl{B_H(0)}\cap S\subset B_S(0)\), so (ii) gives
\(B_S(t)\cap U=\cl{B_H(t)}\cap S\cap U\).
The right side is finite: \(B_H(t)\) has dimension at most
one, and a locally closed algebraic curve has only finitely many
points in its closure outside its original locus.
Because \(B_S(0)\) is closed in \(S\), a curve component of \(B_S(t)\)
not contained in \(B_S(0)\) would have an infinite, nonempty
Zariski-open intersection with \(U\), which is impossible. Such a
curve component must thus be one of the finitely many curve
components of \(B_S(0)\), and smooth continuation cannot switch
between that finite collection.
\end{proof}

\begin{corollary}\label{cor:deep}
The conclusion of the problem holds if \(\Gamma\) has at least two
edges.
\end{corollary}

\begin{proof}
By Lemma~\ref{lem:deep} the \(i\)-th projection of any curve
component of \(\partial V_t\) is \(Y_0\); by hypothesis it is \(Y_t\).
\end{proof}

\begin{lemma}[Coarse continuation in the given, unlabelled stratum]
\label{lem:isolated}
An isolated point \(p\) of \(B_S(0)\) has a unique local continuous
continuation in \(B_S(t)\), and this continuation is holomorphic.
If the curve components of \(B_S(t)\) are fixed and a continuous
section \(p(t)\in B_S(t)\) starts on a curve component of \(B_S(0)\),
then \(p(t)\in B_S(0)\) for all sufficiently small \(t\).
\end{lemma}

\begin{proof}
The smooth trivialization continues the isolated points as
separate zero-dimensional components; holomorphicity
(Lemma~\ref{lem:regularity}) makes their locally unique sections
holomorphic. There are finitely many such sections. For the second
assertion their central values differ from \(p(0)\), so the given
section stays locally in the fixed curve part.
\end{proof}

\begin{lemma}[Stationarity beyond two nodes]\label{lem:deep-points}
If \(S\) has at least three nodes, then
\(B_S(t)\subseteq B_S(0)\) for all \(t\in\Delta\).
In particular, every isolated point of \(B_S(0)\) has constant
local continuation.
\end{lemma}

\begin{proof}
Smooth one node of \(S\) to obtain a boundary stratum \(H\) with at
least two nodes. Lemma~\ref{lem:deep} implies that the curve part
of \(B_H(t)\) is contained in \(B_H(0)\); the remaining components
are finitely many points of \(H\), whose closures do not meet \(S\).
Therefore
\(\cl{B_H(t)}\cap S\subseteq\cl{B_H(0)}\cap S\subseteq B_S(0)\).
Applying (ii) with \(U=\cl{\M}_g\setminus\cl{B_S(0)}\), as in
Lemma~\ref{lem:deep}, gives \(B_S(t)\cap U=\varnothing\); since
\(B_S(0)\) is closed in \(S\), this proves the inclusion.
A continuous section through an isolated point of \(B_S(0)\)
must then be constant near \(0\).
\end{proof}

\begin{proposition}[A stationary corner suffices]\label{prop:fixed-cusp}
Let \(S\) lie on the boundary of the chosen stratum
\(\Delta_\Gamma\). Suppose
\(p\in\cl{B_{\Delta_\Gamma}(0)}\cap S\)
is isolated in \(B_S(0)\) and belongs to \(B_S(t)\) for every
sufficiently small \(t\). Then \(Y_t=Y_0\) for all \(t\). In
particular the conclusion of the problem holds if
\(\cl{B_{\Delta_\Gamma}(0)}\) contains a point with at least three
nodes that is isolated in its full stratum intersection \(B_S(0)\).
\end{proposition}

\begin{proof}
Choose an open neighbourhood \(U\) of \(p\) with
\(B_S(0)\cap U=\{p\}\). Since \(p\) is a limit from the chosen
stratum, the antecedent of (ii) holds for
\(\Delta_\Gamma,S,U\). Thus
\(p\in B_S(t)\cap U=\cl{B_{\Delta_\Gamma}(t)}\cap S\cap U\).
The compactified clutching map from
\(\prod_j\cl{\M}_{g_j,n_j}\) to \(\cl{\M}_g\) is proper. For a fixed \(t\),
the image of its closed subset with \(i\)-th coordinate \(Y_t\)
therefore contains \(\cl{B_{\Delta_\Gamma}(t)}\).
Hence the stable curve \(p\) admits a clutching presentation in
which the \(i\)-th factor is the smooth pointed curve \(Y_t\).
A fixed stable curve has only finitely many such smooth pointed
factors (a normalization component with a labelling of its node
preimages), so \(Y_t\) takes values in a fixed finite set.
Continuity fixes its value near \(0\), and holomorphicity
(Lemma~\ref{lem:regularity}) extends the equality to the whole disc.
For the last assertion combine Lemma~\ref{lem:deep-points} with the
continuation in (i).
\end{proof}

These arguments require isolation in the \emph{whole} stratum
intersection \(B_S(0)\), not just among the limits of the chosen
curve; the next lemma, about a single surface, handles symmetric
two-node corners without any lifted family.

\begin{lemma}[Both partial smoothings at an isolated two-node point]
\label{lem:plumbing}
Let \(Z\subset\M_g\) be an algebraic surface, and let
\(p\in\cl Z\) represent a stable curve with exactly two nodes \(a,b\).
Suppose \(p\) is isolated in the intersection of \(\cl Z\) with its
stratum. In a local deformation chart with the nodes labelled, the
closure of \(Z\) contains a curve approaching \(p\) on which exactly
\(b\) remains nodal, and also such a curve on which exactly
\(a\) remains nodal.
\end{lemma}

\begin{proof}
A versal deformation chart of the stable curve is smooth,
with plumbing coordinates \(u_a,u_b\) (one for each node, vanishing
exactly when that node persists). Its map to the coarse space
is finite and its boundary is \(u_a u_b=0\)
\cite[Sections 8--10, Proposition 12.2 and Corollary 12.3]{HubbardKoch}
(a chart for the single fixed surface \(\cl Z\), not a pullback of a
trivialized family). Choose a lift of \(p\) and an irreducible surface branch \(A\)
of the closure of the lifted interior through that lift; such a
branch exists by finiteness of the chart, and its interior is dense.
Neither \(u_a\) nor \(u_b\) is identically zero on \(A\), so
the zero set of either function on \(A\) has a one-dimensional
component through the chosen point, by Krull's principal ideal
theorem. But \(A\cap\{u_a=u_b=0\}\)
is zero-dimensional there, since the chart is finite and \(p\) is
isolated in its stratum intersection. A curve component of
\(A\cap\{u_b=0\}\) is consequently not contained in \(\{u_a=0\}\);
its general points give the first assertion, and interchanging
\(a,b\) gives the second.
\end{proof}

\begin{corollary}[Equal-genus ends]\label{cor:symmetric}
Suppose the chosen one-node stratum is
\(\M_{h,1}\times\M_{h+r,1}\), with \(r\geq1\), and its whole
intersection with \(\cl{V_t}\) has first coordinate \(Y_t\).
Suppose it has a two-node limit, isolated in the
\emph{whole} two-node stratum intersection of \(\cl{V_t}\),
whose normalization is a chain
\(Y_t\cup_a A\cup_b W\) with \(g(A)=r\), \(g(W)=h\),
where \(W\), including its attaching point, is fixed.
Then \(Y_t=W\).
\end{corollary}

\begin{proof}
Apply Lemma~\ref{lem:plumbing} to \(V_t\), smoothing \(a\) and keeping
\(b\): this produces a curve in the same one-node stratum whose
genus-\(h\) factor tends to \(W\); that factor is identically \(Y_t\)
on the entire stratum intersection, so \(Y_t=W\). Since the two
factors have different genera, no labels on the coarse stratum are
needed.
\end{proof}

\section{Teichm\"uller-geometric input on the central surface}\label{sec:teich}

From now on \(N\subset\T_g\) denotes a marked lift of \(V_0\) to
Teichm\"uller space, a totally geodesic complex surface, and
\(\cl N\) its closure in the bordification of \(\T_g\) used in
\cite{AHW}. For
\(X\in N\), the Teichm\"uller discs through \(X\) tangent to \(N\)
are generated by a two-dimensional space \(Q_XN\) of holomorphic
quadratic differentials on \(X\); \(QN\) is the resulting rank-two
bundle and \(Q_{\max}N\) the locus of differentials in \(QN\) lying
in the largest (generic) stratum of quadratic differentials (by
zero orders, not a boundary stratum). Arana-Herrera and Wright
\cite{AHW} define an equivalence relation on cylinder curves (core
curves of flat cylinders) of
differentials in \(QN\): roughly, two parallel cylinders are
equivalent if they remain parallel, with constant circumference
ratio, on all nearby surfaces of \(QN\), so that they can only be
deformed together. We use the following consequences of their work.
(1) A marked boundary intersection of \(\cl N\) with a stratum is
smooth, irreducible, algebraic, and \emph{semisimple}: a product of
Teichm\"uller-space factors on which the components of the pinched
surface either stay fixed or move along totally geodesic
submanifolds \cite[Theorem 1.3]{AHW}. A one-dimensional such locus
has one positive-dimensional simple factor, and its \emph{active}
(nonconstant) coordinate maps are Teichm\"uller-disc isometries.
(2) Pinched curves occur in whole equivalence classes; each
class has a positive multitwist stabilizing \(N\), and a class can
be unpinched while the other nodes remain pinched
\cite[Corollary 2.2, Lemmas 2.6 and 2.13, Corollaries 2.12 and
3.17]{AHW} (integrality of the twist exponents comes from the
rationality of modulus ratios in Corollary 2.2).
(3) If a nonpoint simple boundary factor is supported on
subsurfaces \(U_1,\ldots,U_s\), then
\(\partial(\bigcup_j U_j)\) is a union of whole equivalence classes
\cite[Lemma 6.1(3)]{AHW}.
Successively unpinching \(r\) distinct
classes gives the \emph{dimension bound} \(\dim L\leq\dim N-r\) for
a boundary locus \(L\) with those classes pinched; see the proof of
\cite[Theorem 3.9]{AHW}. Thus all nodes of a boundary \emph{curve}
of a totally geodesic \emph{surface} form a single equivalence class.

\paragraph{Stabilizers and cusps.}
Let \(L\) be a marked semisimple curve, \(L_j\) an active
projection, and \(\Lambda,\Lambda_j\) their setwise stabilizers in
the product mapping class group and the \(j\)-th mapping class group
(pinching twists already removed), divided by their pointwise
kernels. The disc isometry \(L\to L_j\) identifies \(\Lambda\) with
a subgroup of \(\Lambda_j\) (an element acting trivially on the
active disc acts trivially on \(L\)); the normalized moduli map is
the orbifold covering of this inclusion, of finite degree by
algebraicity, and its source has finite hyperbolic area (a
torsion-free finite cover is an ordinary finite-area curve).
We need a cusp only for the chosen active curve:
Theorem~\ref{thm:zero} makes its differential holomorphic and its
sole mark a zero; forgetting that zero, and taking the orientation
double cover if necessary, changes the lattice only by finite index,
so the noncompactness theorem for algebraic Teichm\"uller curves
\cite{Veech}, \cite[Proposition 2.4]{McMullenSurvey} supplies a cusp.

\begin{lemma}[Forbidden leaf support]\label{lem:forbidden}
A totally geodesic surface cannot have a boundary curve in a stratum
with at least two nodes whose only active normalization component
corresponds to a vertex with exactly one incident half-edge and no
loop.
\end{lemma}

\begin{proof}
Let \(a\) be the unique node incident to that component \(Y\).
At a marked lift of a generic point of the proposed curve, all the
nodes belong to one class by the dimension bound. On the other
hand, the support assertion (3) makes \(\partial Y=\{a\}\) a union of
entire classes. The other node or nodes give a contradiction.
\end{proof}

\paragraph{An arithmetic reduction.}
Let \(\mathcal A_N\) be the affine invariant manifold of square roots
of the largest stratum of \(QN\) (an orbit closure of translation
surfaces, locally linear in period coordinates), and let \(K_N\) be
its field of definition in period coordinates. We use a connected
component and make finite marking choices when needed. For a surface
\(N\), \(\mathcal A_N\) has complex dimension four, rank two, and no
relative tangent directions \cite[Theorem 1.3 and its proof]{WrightTG}.

\begin{proposition}[A singleton forces arithmeticity]\label{prop:field}
If the closure of \(N\) contains a stable Riemann surface with
exactly one node, then \(K_N=\mathbb Q\).
\end{proposition}
\begin{proof}
The node \(a\) is a singleton cylinder class. Realize it
on a differential in \(Q_{\max}N\).
If that differential is a square, its cylinder is an entire
class on the translation surface, and
\cite[Theorem 7.1]{WrightCylinder} (the field of definition lies in
the field generated by circumference ratios within a class) gives
field \(\mathbb Q\).
Otherwise take the orientation cover. Its involution cannot
preserve a flat cylinder: in cylinder coordinates its
derivative is \(-1\), and preservation would give a fixed
regular point. Thus cylinders upstairs come in pairs and
project to cylinders downstairs. The two lifts of \(a\)
have equal circumferences and are equivalent. Any cylinder
equivalent to them would project to an \(N\)-equivalent
cylinder, since the square-root locus locally covers
\(Q_{\max}N\). There are no others. The same theorem, with the
single ratio \(1\), now gives \(K_N=\mathbb Q\).
\end{proof}

\begin{corollary}\label{cor:nonarithmetic}
The conclusion of the problem holds whenever \(K_N\ne\mathbb Q\).
Consequently the only unresolved one-node case has \(K_N=\mathbb Q\).
\end{corollary}
\begin{proof}
A boundary point has a compatible marked lift: use a local branch
of the Dehn space as in the proof of \cite[Lemma 3.13]{AHW}.
Proposition~\ref{prop:field} rules out a one-node stratum when
\(K_N\ne\mathbb Q\). The chosen stratum must then have at least
two nodes, and Corollary~\ref{cor:deep} applies.
\end{proof}

\section{The active differential vanishes at the attachment}\label{sec:sewing}

Only the central totally geodesic surface is used in this section.
By Corollaries~\ref{cor:deep} and \ref{cor:nonarithmetic}
the remaining case is a one-node stratum; a nonseparating one-node
stratum has a single factor, which cannot be constant on a curve, so
the node is separating. Let \(N\subset\T_g\) be a marked lift of
\(V_0\) with one-node boundary curve
\(L=\{(Y,y)\}\times D\), \(D\subset\T_{m,1}\),
where \((Y,y)\) is the fixed pointed curve of genus \(h\ge1\),
\(D\) is a Teichm\"uller disc, and \(a\) is the separating node;
by the dimension bound the class of \(a\) is a singleton. Locally on
\(D\), use a parameter \(\tau\), write the pointed curve at \(\tau\)
as \((E_\tau,x_\tau)\), and choose a nonzero holomorphic generator
\(q_D(\tau)\) of its Teichm\"uller differential line, a meromorphic
quadratic differential on \(E_\tau\) with at most a simple pole at
\(x_\tau\). We recall some terminology.
A quadratic differential with a double pole has a \emph{quadratic
residue} \(R^2\), the coefficient of \((dz/z)^2\); if \(R^2\ne0\)
there is a local \emph{exponential coordinate} \(w\), unique up to a
scalar, in which it equals \(R^2(dw/w)^2\). Given two such
differentials (``caps'') with matching residues on the two sides of
a node, identifying \(w_Yw_E=u\) produces a smooth curve with a flat
cylinder replacing the node: \emph{sewing} with parameter \(u\).
The \emph{quadratic-Hodge bundle} \(\mathcal H\) over the
deformation space consists of quadratic differentials with at most
double poles and matching residues at the node. A
\emph{Jenkins--Strebel} (or Strebel) differential has closed
horizontal trajectories off a critical graph, sweeping out flat
cylinders; it is \emph{complete} for a multicurve if its cylinder
cores are exactly that multicurve.

\begin{lemma}[A boundary pencil and its sewings]\label{lem:cap-pencil}
Locally at every point of \(D\) there are meromorphic quadratic
differentials
\(Q(\tau,s)=\bigl(q_Y,\ q_E^0(\tau)+s q_D(\tau)\bigr)\)
on \(Y\sqcup E_\tau\), depending holomorphically on \((\tau,s)\)
(\(q_E^0(\tau)\) is identified in the proof),
with matching nonzero quadratic residue
\(R^2\), independent of \(\tau,s\), at the node.
They belong to the quadratic-Hodge closure of \(QN\).
Sewing their double poles in exponential coordinates
\(w_Yw_E=u\) gives a curve in \(N\), for every sufficiently
small \(u\ne0\). The cap \(q_Y\), and hence \(w_Y\), is
independent of both \(\tau\) and \(s\).
\end{lemma}
\begin{proof}
Work in a labelled finite deformation chart of the central
locus. Close the graph of \(X\mapsto Q_XN\subset\mathcal H_X\),
\(X\in N\), in the Grassmannian bundle of two-planes in
\(\mathcal H\). For any limiting plane \(W\) at \(P=(Y,E_\tau)\),
its residue kernel (the differentials in \(W\) with zero residue at
the node) is contained in \(\mathbb C(0,q_D(\tau))\), by
\cite[Corollaries 3.11--3.12 and the proof of Theorem 3.9]{AHW};
by dimension it equals this line.

\emph{Uniqueness of the limiting plane} (not merely generic
finiteness). For \(X_n\to P\) realizing \(W\), take the complete
one-cylinder
Jenkins--Strebel differentials \(J_n\in Q_{X_n}N\)
given by \cite[Corollary 2.7]{AHW}, with core \(a\)
and circumference one.
They are bounded in a local trivialization of \(\mathcal H\):
otherwise dividing by their divergent norms gives a nonzero
limiting differential but horizontal circumference tending
to zero; a small natural-coordinate disc at a regular
point of that limit contains a uniform smaller disc for
the approximants, and a regular point off their critical graphs
would lie on a closed horizontal trajectory too short
to leave this disc, impossible in a flat coordinate.
For any convergent subsequence of \(J_n\), cylinder/residue
convergence, as in \cite[Lemma 3.10]{AHW}, gives residue
\(-1/(4\pi^2)\) at the node. Almost every limiting horizontal
trajectory is closed: away from the critical trajectories,
flows converge for their common period one, and a bounded
length trajectory cannot escape down an infinite pole cylinder.
Each resulting closed leaf is peripheral to the node, since
the approximating leaves have core \(a\).
Strebel uniqueness \cite[Theorem 2]{Zvonkine} identifies the
limit as \(J(P)=(J_Y,J_{E_\tau})\),
the pair of one-pole Strebel differentials of perimeter one,
with \(J_Y\) fixed. Thus every \(W\) equals
\(\mathbb C J(P)+\mathbb C(0,q_D(\tau))\).
The closure of the graph of \(\tau\mapsto W(\tau)\) over \(L\)
(with its reduced structure) is now proper with singleton
fibres, hence finite and birational over the smooth curve
\(L\), and therefore isomorphic to it. Thus \(W(\tau)\) is
holomorphic even at formerly exceptional points.
A holomorphic fixed-residue section of \(W(\tau)\) has
\(Y\)-part equal to a fixed multiple of \(J_Y\); adding
\(s(0,q_D)\) gives the asserted pencil.

\emph{Sewing an arbitrary vector of the pencil} (not only a
differential selected by an existence theorem).
Approximate the given vector by \((X_n,Q_n)\in Q_{\max}N\). Compact annuli near
both limiting double poles are zero-free. On its approximating
annulus the nonzero core period \(L_n=\oint_a\sqrt{Q_n}\) tends
to \(2\pi iR\), up to sign, and the coordinate
\(w_n=\exp\bigl(\frac{2\pi i}{L_n}\int\sqrt{Q_n}\bigr)\)
is single-valued. Choose its multiplicative constant so it converges
to the limiting exponential coordinate on a slightly larger
annulus. On a smaller annulus it is then univalent, by
uniform convergence and the argument principle, and its
image contains a common round subannulus. After rotating
\(Q_n\), these round circles are closed horizontal
trajectories, so \(a\) is a cylinder of \(Q_n\).
The two annuli belong to the same maximal \(a\)-cylinder:
an annulus between homotopic flat cores on a holomorphic
quadratic differential contains no zeros, by Gauss--Bonnet.
Cutting this cylinder and filling both sides by punctured
discs in the displayed coordinates is holomorphic locally, and
convergence on compact subsets on both sides gives convergence
of these caps to the prescribed pair. The singleton
Cylinder Deformation Theorem \cite[Theorem 5.1]{WrightCylinder}
allows arbitrary shear and positive height for this cylinder
without changing the caps, while staying in \(\mathcal A_N\).
For each fixed sufficiently small \(u\ne0\) (uniformly over a
relatively compact cap-parameter neighbourhood: common
pole-coordinate radii keep the cylinder height positive), choose
these deformations to have sewing parameter \(u\); changing
coordinate scalars only reparametrizes \(u\). The sewn approximants
converge to the sewing of \(Q(\tau,s)\), which lies in \(N\) by
closedness. Integer changes of the sewing marking are harmless
because the Dehn twist \(T_a\) stabilizes \(N\); the marked lift of
the holomorphic sewing family is taken on a simply connected sector
of the punctured \(u\)-disc, winding around zero giving that twist.
No simultaneous realization of intersecting cylinder multicurves is
required.
\end{proof}

The next lemma is the heart of the obstruction: sewing the pencil
with parameter \(u\) yields \emph{three} parameters \((\tau,u,s)\)
unless the active differential vanishes at the attaching point.

\begin{lemma}[Second-order sewing obstruction]\label{lem:sewing}
For the pencil in Lemma~\ref{lem:cap-pencil}, suppose
the unpointed curves \(E_\tau\) vary.
If \(q_D\) has either a simple pole or a regular nonzero
value at \(x_\tau\), the map \((\tau,u,s)\mapsto\M_g\) given by the
sewn curves has complex rank three at suitable points with
arbitrarily small \(u\ne0\).
\end{lemma}
\begin{proof}
Choose a coordinate \(z\) at \(x_\tau\), holomorphically
in \(\tau\), for which \(q_E^0=R^2(dz/z)^2\), and write
\(q_D=(d(\tau)z^{-1}+e(\tau)+O(z))\,dz^2\),
\(A=sd/R^2\), \(B=se/R^2\).
The exponential coordinate of \(q_E^0+s q_D\), normalized
by \(w_E/z\to1\), satisfies
\(w_E=z+\frac A2z^2+(\frac{A^2}{16}+\frac B4)z^3+O(z^4)\).
For the \(A\)-normalized fundamental bidifferential write
\(\mathcal B(w,w')=((w-w')^{-2}+\beta+O(w,w'))\,dw\,dw'\).
A change \(w=z+az^2+bz^3+\cdots\) changes its constant
regular coefficient by \(a^2-b\): substitute
\(z=w-aw^2+(2a^2-b)w^3+\cdots\) in the singular part;
the holomorphic part has unchanged constant coefficient.
Consequently
\(\beta_E(\tau,s)=\beta_E(\tau,0)+3s^2d(\tau)^2/(16R^4)
-se(\tau)/(4R^2)\).
If \(d\ne0\), its \(s\)-derivative is not identically
zero; if \(d=0,e\ne0\), that derivative is everywhere
nonzero.

Let \(\eta_\alpha,\omega_j\) be normalized holomorphic
one-forms on \(E_\tau,Y\), and set
\(e_\alpha=(\eta_\alpha/dw_E)(x_\tau)\),
\(v_j=(\omega_j/dw_Y)(y)\), \(c_*=2\pi i\).
These evaluations are independent of \(s\), because
\(dw_E/dz(0)=1\) and the \(Y\)-cap is independent of \(s\).
Write \(\Omega\) for the sewn period matrix and
\(\Omega_E,\Omega_Y\) for those of its components.
The separating-node expansion is
\begin{align}
 \Omega_{\alpha\beta}
    &=(\Omega_E)_{\alpha\beta}
       +c_*u^2\beta_Y e_\alpha e_\beta+O(u^3),\nonumber\\
 \Omega_{\alpha j}
    &=-c_*u e_\alpha v_j+O(u^2),\label{eq:general-periods}\\
 \Omega_{ij}
    &=(\Omega_Y)_{ij}
       +c_*u^2\beta_E v_i v_j+O(u^3).\nonumber
\end{align}
Here the \(E\)-block's displayed coefficients are independent
of \(s\), and \(\Omega_Y\) is constant.
(Continuation of \(\omega_j\) to \(E\) starts with
\(-uv_j\,dw_E/w_E^2\), whose normalized extension has
constant regular coefficient \(-uv_j\beta_E\); returning to \(Y\)
gives \(u^2v_j\beta_E\,dw_Y/w_Y^2\), with \(B_i\)-period
\(c_*u^2\beta_Ev_iv_j\). Reversing the components gives the
\(E\)-block.) The analytic matching construction
\cite[Theorem 3.3, Section 5.1.1]{HuNorton} controls the
remainders; their \(\tau,s\)-derivatives have the same
orders, and a \(u\)-derivative lowers the order by one.

By Torelli's theorem, some entry
\(F(\tau)=(\Omega_E)_{\rho\sigma}\) is nonconstant.
Choose a general \(\tau\) with \(F'(\tau)\ne0\), and
choose \(\alpha,j\) with \(e_\alpha v_j\ne0\), using
base-point freeness of the canonical systems.
At an \(s\) with \(\partial_s\beta_E\ne0\),
\(\det\,\partial(\Omega_{\rho\sigma},\Omega_{\alpha j},\Omega_{jj})
/\partial(\tau,u,s)
=-c_*^2F'e_\alpha v_j^3(\partial_s\beta_E)u^2+O(u^3)\ne0\).
The period matrix is holomorphic in the marked curve, so the
moduli rank is at least three.
\end{proof}

\begin{theorem}[A positive-order zero is necessary]\label{thm:zero}
If a totally geodesic surface has a one-node separating boundary
curve with one constant factor, then the active Teichm\"uller
differential vanishes at the attaching point, and the active
component has genus \(m\ge2\).
\end{theorem}
\begin{proof}
First, the unpointed active curve is nonconstant: for \(m=1\),
forgetting the sole marking changes no elliptic modulus, and for
\(m\ge2\), a constant unpointed projection would put the marked Teichm\"uller disc in
a Bers fibre (a fibre of the forgetful map
\(\T_{m,1}\to\T_m\), \cite[Section 9(B)]{BersFinite}). Kra's strict
comparison \cite[Theorem 3, Section 10]{Kra} gives, for distinct
points in that fibre, ambient Teichm\"uller distance
strictly less than intrinsic Poincar\'e distance, so
Schwarz's lemma for a disc mapping into the fibre
contradicts the ambient disc isometry.
A Teichm\"uller differential has finite area, hence at most a
simple pole at the sole marking. Unless it vanishes there,
Lemmas~\ref{lem:cap-pencil} and \ref{lem:sewing} give
a rank-three family in the two-dimensional \(N\),
a contradiction. The singularity order at the marked
point is constant along its generating disc, so it
suffices to work at the general points used above.
Finally, on an elliptic curve,
\(H^0(K^2(p))=H^0(\mathcal O(p))\) is one-dimensional, and
every nonzero member is a nowhere-zero holomorphic
quadratic differential; thus \(m=1\) is impossible.
\end{proof}

\begin{proposition}[Arithmeticity of the active curve]\label{prop:active-field}
In the configuration of Theorem~\ref{thm:zero}, the square-root
orbit of the active differential has field of definition
\(\mathbb Q\); in particular it contains a square-tiled surface.
\end{proposition}
\begin{proof}
The differential \(q_D\) is holomorphic, so forgetting the
attachment preserves its Teichm\"uller disc. Its unpointed
image is an algebraic Teichm\"uller curve (the forgetful map is
proper, and the marked zero has finitely many choices).
Choose a periodic direction of \(q_D\), with cylinder cores
\(B=\{b_1,\ldots,b_r\}\). Its cusp gives a point of \(\cl N\)
with nodes \(a\cup B\); no core is contracted after forgetting the
mark, since a holomorphic quadratic differential has neither an
inessential flat core nor two homotopic distinct maximal cylinders,
by Gauss--Bonnet. Since \(\{a\}\) is one class, the dimension bound
shows that \(B\) is exactly one further \(N\)-class.

Proposition~\ref{prop:field} gives \(K_N=\mathbb Q\).
On its rational square-root tangent space, integrations over
the lifted cylinder cores are rational linear functionals
\(\lambda_j\), and equivalence makes
\(\lambda_j=c_j\lambda_1\) with \(\lambda_1\ne0\).
Evaluating on a rational tangent vector where \(\lambda_1\ne0\)
gives \(c_j\in\mathbb Q\) \cite[Lemma 7.2]{WrightCylinder}.
The finite-area vector \((0,q_D)\) belongs to \(\cl{QN}\), and its
cylinders persist on approximants in \(Q_{\max}N\): the normalized
annular coordinates in the proof of Lemma~\ref{lem:cap-pencil}
apply to compact subcylinders. By \cite[Corollary 2.2]{AHW}, these
\(c_j\) are the same constants on every approximant with those
marked cylinders, depending only on the cylinder curves and \(N\).
Cylinder circumferences converge to those of \(q_D\), so its
ratios equal these fixed rational numbers (not merely an arbitrary
limit of rational ratios).

Take the orientation cover of this \emph{smooth} active
differential, or the differential itself if it is a square.
In the nonsquare case every downstairs cylinder has two
equal-circumference lifts, and these exhaust the upstairs
cylinders: all zeros remain zeros, and a cylinder invariant under
the orientation involution would contain a fixed regular point.
Thus all horizontal circumference ratios upstairs are rational, and
\cite[Theorem 1.9]{WrightCylinder} gives field \(\mathbb Q\). The
closed square-root orbit has complex dimension two; rational real
and imaginary relative periods are dense in its rational period
plane, and after scaling one such point has integral periods and is
square-tiled. No specialization of a possibly disconnected
orientation cover is used.
\end{proof}

The corner argument needs isolation of the two-node cusps of \(L\)
(points of \(\cl L\) where one further curve \(b\) of \(E_\tau\) is
pinched) in the \emph{whole} two-node stratum intersection \(B_S(0)\).

\begin{proposition}[Isolation of two-node cusps]\label{prop:two-corner}
Suppose a cusp of a central one-node boundary curve has exactly
two nodes. Then its coarse point \(p\) is isolated in the
whole intersection \(B_S(0)\) with its two-node stratum \(S\).
\end{proposition}
\begin{proof}
Choose a compatible marked lift \(P\in\cl N\), with nodes
\(a,b\). They are two singleton classes, \(a\) being already a
singleton along the one-node curve. For \(X_n\in N\) tending
to \(P\), the complete one-class Strebel differentials
\(J_{a,n},J_{b,n}\), each normalized to circumference one,
have canonical limits \(J_a(P),J_b(P)\): the bounded-trajectory
argument of Lemma~\ref{lem:cap-pencil} applies verbatim. For each,
the other node has zero residue, since a nonzero residue would produce
an additional cylinder on the approximants. On each
normalization component the limit is the Strebel differential
with its prescribed positive pole perimeters, other node
preimages being retained as punctures, and components with no such
pole carry zero; this uses the punctured uniqueness theorem
\cite[Theorem 2, Remark 3.8]{Zvonkine}.
Their residue vectors are independent, so every limiting
plane is \(W(P)=\mathbb C J_a(P)+\mathbb C J_b(P)\).

\emph{All two-cylinder caps lie in \(W(P)\).}
For any \(r>0\), let \(J_{1,r}(P)\) denote the complete
Strebel cap differential with perimeters \(1,r\) at the two
nodes. We claim \(J_{1,r}(P)\in W(P)\).
On each fixed \(X_n\), positive integral multitwists
\(T_a^kT_b^l\) stabilize \(N\), and their horocycle differentials
\cite[Lemmas 2.4--2.5]{AHW} are complete two-cylinder
Strebel differentials of prescribed modulus ratio \(k/l\).
Continuity and density allow all positive ratios, and the
projective endpoints are \(J_{a,n},J_{b,n}\): on a fixed
smooth \(X_n\), the flat lengths of the fixed curves \(a,b\) have
positive lower and finite upper bounds on the compact sphere of
unit-area differentials, so a cylinder whose modulus ratio tends
to zero has height and area tending to zero, its transverse weight
vanishes, the limiting horizontal measured foliation is
supported on the remaining core, and uniqueness gives the
one-class limit.
The circumference ratio varies continuously between these
endpoints; the first endpoint tends to zero as \(n\to\infty\)
and the second to infinity, because a small
circle about a zero-residue node in \(J_a(P)\) or \(J_b(P)\)
has length tending to zero with its radius, and its approximating
circles bound the opposite core length from above.
Thus one may choose two-cylinder differentials on \(X_n\)
with circumferences exactly \(1,r\). They are bounded by the same
short-trajectory argument, and their limits are \(J_{1,r}(P)\) by
Strebel uniqueness, proving the claim. Independent cylinder
deformations, as in Lemma~\ref{lem:cap-pencil}, consequently put
\emph{all} sufficiently small sewings of these caps in \(N\).

\emph{A second branch would be a translate of \(N\).}
Suppose another local surface branch over \(p\) has a curve
in the full two-node stratum. A compatible marked lift
\(N'\) then has \(\{a,b\}\) as one class, with some fixed
circumference ratio \(r\), and its complete Strebel limits are
the same \(J_{1,r}(P)\).
Classwise cylinder deformation and passage to the limit
give a nonconstant sewing curve
\(u_a=c_a z^{d_a}\), \(u_b=c_b z^{d_b}\) with \(c_a c_b\ne0\),
\(d_a,d_b>0\), in this branch; the positive integral exponents come
from the rational modulus ratio, and the constants are obtained
by a subsequence (in exponential coordinates normalized
at the cap critical graphs their moduli are fixed and
their phases lie in a compact torus).
This curve also belongs to the two-parameter sewing family
in \(N\). On a simply connected sector, after a Dehn-twist
translate of \(N'\), it lifts to a common holomorphic curve
in \(N\cap N'\), tending to the same marked \(P\).

Here smoothness of the \emph{coarse interior} \(V_0\) is
decisive. At every point \(X\) of the common curve, the
finite stabilizer of \(X\) acts transitively on the local
branches over the irreducible smooth germ of \(V_0\), so
some mapping class fixes \(X\) and takes \(N\) to \(N'\).
There are countably many mapping classes, and their fixed sets
on the common parameter curve are analytic; unless one
contains the whole curve, their union is countable.
Hence one such \(\gamma\) fixes the whole curve and takes
\(N\) to \(N'\). Continuity gives \(\gamma P=P\), so it permutes
\(a,b\) and must send their two singleton classes for
\(N\) to two singleton classes for \(N'\), a contradiction.
No curve of \(B_S(0)\) passes through \(p\), proving isolation.
\end{proof}

\begin{corollary}[A continued corner without (L)]\label{cor:corner-section}
At such a two-node cusp, (i)--(ii) give a unique local
holomorphic section \(p(t)\in B_S(t)\), lying in the closures
of both incident one-node stratum intersections.
This statement does not assert that \(p(t)\) is fixed.
\end{corollary}
\begin{proof}
Use Proposition~\ref{prop:two-corner} and
Lemma~\ref{lem:isolated}. Both partial smoothings occur centrally
by Lemma~\ref{lem:plumbing}. Isolation permits (ii) in a
neighbourhood of \(p\), for either one-node stratum.
\end{proof}

\section{Curve rigidity and the lifting property (L)}\label{sec:lifting}

To run the corner mechanism in the one-node case we need to know
how the \emph{deformed} boundary curves \(C_t\subset\partial V_t\)
look; the only tool available is the rigidity of Teichm\"uller
curve maps, which requires maps from curves to moduli spaces with
fixed monodromy. The following condition, stated separately and
\emph{not} inferred from the hypotheses of the problem, supplies
this. (A \emph{fine level cover} of \(\M_{h,n}\) is a finite cover
that is a fine moduli space, so maps into it have well-defined
monodromy.)

\noindent\textbf{(L)} \emph{The continued normalization of each
boundary curve used below admits a holomorphic family of connected
finite algebraic covers with smoothly trivial source, mapping
holomorphically to fixed fine level covers of the labelled moduli
factors. The images are the given continued boundary curves. On the
central fibre these maps are finite orbifold covers of the marked
semisimple curve maps.}

\noindent Only the curve \(C_t\) of Theorem~\ref{thm:conditional}
and, when used, its second curve \(C'_t\) need satisfy (L); isolated
points need no lifting.

\begin{lemma}[Pointed curve-map rigidity]\label{lem:pointed}
Let \(f_t:C_t\to\M_{h,n}\) be a holomorphic family of maps from
smoothly trivial finite-area algebraic curves. Assume that the maps
are locally liftable to \(\T_{h,n}\) and that their
mapping-class-group monodromy representations \(\rho\) are the same
under the given topological identification of the sources (as holds
for an actual family of maps to a fixed fine level
cover with smoothly trivial source). If \(f_0\) is a
Teichm\"uller-disc quotient map, then the pairs \((C_t,f_t)\) are
isomorphic to \((C_0,f_0)\), with the prescribed marking of the source.
\end{lemma}

\begin{proof}
\review{There was no need to adapt McMullen's argument---it carries over verbatim.}

This is McMullen's argument \cite{McMullen}, spelled out to
cover the pointed version. Normalize the hyperbolic metrics
to curvature \(-4\), so that a Teichm\"uller disc is isometric
for \(d_{\T}=\frac12\log K\). A hyperbolic element \(\gamma\) of
the central source group acts on the target by a pseudo-Anosov
mapping class, whose Teichm\"uller translation length is precisely
the length \(\ell_0(\gamma)\) of the corresponding source geodesic.
Schwarz's lemma for a locally liftable map (the Teichm\"uller
metric is the Kobayashi metric, by Royden's theorem
\cite[Section 7(A)]{BersFinite}) gives
\(\ell_t(\gamma)\geq\inf_{x\in\T_{h,n}}d(x,\rho(\gamma)x)
=\ell_0(\gamma)\); the fixed representation \(\rho\) is essential
here. The non-domination theorem for marked length spectra of
finite-area hyperbolic surfaces, used in \cite{McMullen}, now
identifies \(C_t\) with \(C_0\) as marked Riemann surfaces.

Under this identification lift the maps to the disc and
\(\T_{h,n}\), with the same equivariance. On the axis of each
hyperbolic \(\gamma\), equality in the inequality forces
the lift to map isometrically onto the unique Teichm\"uller axis
of \(\rho(\gamma)\); uniqueness and realization of this axis hold
for punctured surfaces as well
\cite[Theorem 2.1 and its following remark]{Tanigawa}.
Choose two source axes crossing transversely. Their two fixed
target axes meet at the central image of the intersection, and
at no other point. The lift therefore agrees with the central
lift at that point and along either axis; equivariance fixes its
orientation along the axis. The identity theorem gives equality
of the holomorphic lifts everywhere.
\end{proof}

\begin{lemma}[Product form under (L)]\label{lem:curve}
For a boundary-curve family satisfying (L), whose central map
\(f=(f_j)_j\) to the labelled factors is semisimple, the continued
maps \(F_t=((F_t)_j)_j\) have the following form after passing
to the covers in (L) and identifying each source with the central
curve \(C\):
\begin{equation}\label{eq:product}
 (F_t)_j(z)=
 \begin{cases}
 f_j(z),&j\in I,\\
 Z_j(t),&j\notin I,
 \end{cases}
\end{equation}
where \(I\) is the set of initially active coordinates and the
\(Z_j\) are holomorphic. Every central cusp consequently gives a
holomorphic family of stable limits of the same dual graph; all
initially active stable factor limits are fixed.
\end{lemma}

\begin{proof}
Work on the connected finite covers supplied by (L). For an
initially constant coordinate the map to its fine level space is
constant on the central fibre; its monodromy is trivial and
remains so by smooth triviality of this actual family of level
maps, so on every fibre it lifts to \(\T_{h_j,n_j}\).
Bers' bounded-domain realization
\cite[Section 6]{BersFinite} shows that this lift
is constant: its bounded coordinate functions extend across the
finitely many punctures of the algebraic source, and hence are
constant on its compactification. Evaluating on a local
holomorphic section of the source family shows that
\(Z_j(t)\) is holomorphic.
Apply Lemma~\ref{lem:pointed} to any one active coordinate to
identify the marked source with \(C\), and its axis argument
to the other active coordinates with the same marked
identification; all are then the fixed maps \(f_j\). This proves
\eqref{eq:product}.

Let \(c\) be a puncture of the compactification of \(C\).
The fixed algebraic maps \(f_j\) extend to the proper spaces
\(\cl{\M}_{h_j,n_j}\), and their values at \(c\) are fixed stable
pointed curves. The other factors are the smooth curves \(Z_j(t)\).
Gluing their prescribed marked points according to the original
stratum gives the asserted holomorphic stable-curve family
in the coarse compactification. Fixed stable factors do not
smooth or acquire nodes, and the remaining factors stay smooth,
so its graph is constant. No lifted cusp family was assumed.
\end{proof}

\section{Conditional completion by coarse corner comparison}\label{sec:conditional}

\begin{theorem}\label{thm:conditional}
Under the hypotheses of the problem and the additional
boundary-curve property (L), \(Y_t=Y_0\) for all \(t\).
\end{theorem}

\begin{proof}
Work near \(t=0\). By Corollary~\ref{cor:deep} and the reduction
opening Section~\ref{sec:sewing}, write
\(\Delta_\Gamma=\M_{h,1}\times\M_{m,1}\), with first
coordinate \(Y_t\) and separating node \(a\). Choose a continued curve
component \(C_t\) of \(\partial V_t\). Theorem~\ref{thm:zero} gives
\(m\ge2\), and (L) with Lemma~\ref{lem:curve} gives a fixed
Teichm\"uller curve \(D\subset\M_{m,1}\) (the image of the disc
\(D\) of Section~\ref{sec:sewing}) such that
\begin{equation}\label{eq:Ct}
 C_t=\{Y_t\}\times D.
\end{equation}
Suppose \(Y_t\) is nonconstant. A cusp of \(D\) has a stable
pointed limit \(R\), with nonempty node set \(B\). Clutching
gives a holomorphic section
\begin{equation}\label{eq:p}
 p(t)=Y_t\underset a\cup R\in\cl{C_t}.
\end{equation}
Its stratum \(S\) has nodes \(\{a\}\cup B\).

\emph{Step 1: \(p(0)\) is isolated in \(B_S(0)\).}
For \(|B|=1\), Proposition~\ref{prop:two-corner} gives
isolation of \(p(0)\). For \(|B|\ge2\), if \(p(0)\) belonged
to a curve component of \(B_S(0)\),
Lemmas~\ref{lem:deep} and \ref{lem:isolated} would put
\(p(t)\) in the fixed central curve part. This arc is
nonconstant, since a fixed stable curve has only finitely
many smooth pointed normalization components. In the finite
product cover labelling those components, the explicit arc
\eqref{eq:p} varies only in its \(Y\)-coordinate; a branch lies
in a central algebraic curve, and the identity theorem makes its
other coordinates constant. Its sole active component would be the
leaf \(Y\), contrary to Lemma~\ref{lem:forbidden}.
Thus \(p(0)\) is isolated in the whole \(B_S(0)\), and its
continuation \(p(t)\) is unique and isolated for small
\(t\). If \(|B|\ge2\), the last assertion of
Proposition~\ref{prop:fixed-cusp} finishes. Hence assume \(B=\{b\}\).

\emph{Step 2: the second boundary curve.}
Choose compatible marked lifts of \(C_0,p(0)\) in \(\cl N\).
Unpinching the singleton class \(a\), while retaining \(b\),
produces a boundary curve \(C'_0\) approaching \(p(0)\):
the marked intersection is smooth and irreducible,
contains points approaching a point outside its stratum,
and has dimension at most one. Let \(U\) be the normalization
component obtained by smoothing \(a\). Its coordinate on
\(C'_0\) is active: a positive power of \(T_a\) stabilizes
this marked boundary curve and acts on \(\T(U)\) as an
infinite-order Dehn twist, which cannot fix a point
(\(a\) remains essential in \(U\), being essential and
not isotopic to \(b\) on the original surface).

\emph{Step 3a: \(b\) nonseparating.}
Then \(U\) is the only coordinate of \(C'_0\). Apply (L) and
Lemma~\ref{lem:curve} to its continuation \(C'_t\). The chosen
cusp gives the constant section \(p(0)\in B_S(t)\). Uniqueness of
the continuation (Lemma~\ref{lem:isolated}) identifies it with
\(p(t)\), fixing \(Y_t\) by Proposition~\ref{prop:fixed-cusp}.

\emph{Step 3b: \(b\) separating.}
The section has the chain presentation
\(Y_t\cup_a A\cup_b W\) with \(g(A)=r\), \(g(W)=s\),
where \(A,W\), including their attachments, are fixed (pieces of
the fixed \(R\) of \eqref{eq:p}, hence still resting on (L)).
Stability gives \(r,s\ge1\): a rational component of the
once-pointed one-node \(R\) would have at most two special
points. If \(s=h\), Corollary~\ref{cor:symmetric} directly
gives \(Y_t=W=Y_0\); this uses constancy on the \emph{whole}
one-node stratum intersection and requires no choice of
end labels. If \(s\ne h\), apply (L) and Lemma~\ref{lem:curve} to
\(C'_t\). Its chosen cusp yields \(q(t)\in B_S(t)\), \(q(0)=p(0)\),
whose stable \(U\)-coordinate is the fixed curve \(Y_0\cup_a A\).
Uniqueness of the continuation gives \(q(t)=p(t)\). The genus-\(h\) end
is intrinsically distinguished, since the end genera differ;
its pointed normalization component is respectively \(Y_0\)
and \(Y_t\). Thus \(Y_t=Y_0\) in this case too.
Local constancy extends to the disc by Lemma~\ref{lem:regularity}.
\end{proof}

\section{Remaining open issues}\label{sec:open}

The unresolved case is a one-node separating stratum with
\(K_N=\mathbb Q\), an arithmetic active curve of genus \(m\ge2\),
and an attaching zero of order \(k\ge1\).
The canonical limiting plane and fixed cap in
Lemma~\ref{lem:cap-pencil} control \(V_0\), not the normalized
boundary covers of its deformation. At two-node cusps,
Corollary~\ref{cor:corner-section} supplies a continued
corner from (i)--(ii) alone, but its normalization components
can a priori all move; proving stationarity of that corner,
or fixing its active remainder, would close this case.
Theorem~\ref{thm:conditional} uses (L) first to fix the active
pointed image in \eqref{eq:Ct}, and then for the second curve
in the nonseparating and unequal-end-genus cases; for cusps with
more than two nodes, its isolation argument also depends on
\eqref{eq:Ct}. Neither (i)--(ii) nor the arithmetic reduction
supplies the fixed-monodromy orbifold maps that (L) demands; a
source-lifting theorem, or a surface-level replacement, is still
needed. An alternative route would produce a fixed nodal unpointed remainder
of the active component in every deformed closure and compare its
possible attachments; no such persistent remainder has been
obtained from (i)--(ii).

Two warnings (not moduli-space counterexamples) show why more
control is needed: the normalized compact curves
\(z^2=(x^3-1)/(1+tx^2)\) have genus one at \(t=0\) and genus two
nearby; and the graphs \(y=x^2+tx^3\) have total projective central
equation \(Z(YZ-X^2)=0\) but individual central closure
\(YZ-X^2=0\) (products with \(\mathbb P^1\) give surfaces with fixed
reduced boundary). Lemma~\ref{lem:regularity} supplies boundary
holomorphicity, not preservation of finite-cover topology or
exclusion of extra central support.

\begin{thebibliography}{99}

\bibitem{AHW}
F. Arana-Herrera and A. Wright,
\emph{The geometry of totally geodesic subvarieties of moduli spaces
of Riemann surfaces}, arXiv:2412.16330v2, 2025.

\bibitem{BersFinite}
L. Bers,
\emph{Finite dimensional Teichm\"uller spaces and generalizations},
Bull. Amer. Math. Soc. (N.S.) \textbf{5} (1981), 131--172.

\bibitem{HubbardKoch}
J. H. Hubbard and S. Koch,
\emph{An analytic construction of the Deligne--Mumford
compactification of the moduli space of curves},
J. Differential Geom. \textbf{98} (2014), 261--313.

\bibitem{HuNorton}
X. Hu and C. Norton,
\emph{General variational formulas for Abelian differentials},
Int. Math. Res. Not. \textbf{2020}, 3540--3581; arXiv:1706.05366v2.

\bibitem{Kra}
I. Kra,
\emph{On the Nielsen--Thurston--Bers type of some self-maps
of Riemann surfaces},
Acta Math. \textbf{146} (1981), 231--270.

\bibitem{McMullen}
C. T. McMullen,
\emph{Rigidity of Teichm\"uller curves},
Math. Res. Lett. \textbf{16} (2009), 647--649.

\bibitem{McMullenSurvey}
C. T. McMullen,
\emph{Billiards and Teichm\"uller curves},
Bull. Amer. Math. Soc. \textbf{60} (2023), 195--250.

\bibitem{Stacks}
The Stacks Project Authors,
\emph{The Stacks Project}, Tags 037P, 052A, 0559, 0578.

\bibitem{Tanigawa}
H. Tanigawa,
\emph{Orbits and their accumulation points of cyclic subgroups
of modular groups},
Tohoku Math. J. (2) \textbf{43} (1991), 289--299.

\bibitem{Veech}
W. A. Veech,
\emph{Teichm\"uller curves in moduli space, Eisenstein series and
an application to triangular billiards},
Invent. Math. \textbf{97} (1989), 553--583.

\bibitem{WrightCylinder}
A. Wright,
\emph{Cylinder deformations in orbit closures of translation surfaces},
Geom. Topol. \textbf{19} (2015), 413--438.

\bibitem{WrightTG}
A. Wright,
\emph{Totally geodesic submanifolds of Teichm\"uller space},
J. Differential Geom. \textbf{115} (2020), 565--575;
corrected version arXiv:1702.03249v3, 2024.

\bibitem{Zvonkine}
D. Zvonkine,
\emph{Strebel differentials on stable curves and
Kontsevich's proof of Witten's conjecture},
arXiv:math/0209071v2, 2004.

\end{thebibliography}

\end{document}
